\section{Introdução comparativa de Tomografia de Tempo de Trânsito}
\label{sec:travel-time_tomography}

A primeira revisão bibliográfica realizada tratou da \gls{ttt} --- dado que essa técnica é a que melhor se encaixa nas limitações do projeto --- buscando estabelecer suas aplicações, hipóteses de operação, técnicas adjacentes, e também construir um modelo-padrão básico que descreva a estrutura matemática básica de operação dela.

Foram buscados 80 artigos nas plataformas da \textsl{IEEExplore} e \textsl{Google Scholar}, utilizando os termos ``\ingles{Travel-Time Tomography}'' e ``\ingles{TTT}'', dos quais 10 artigos foram selecionados por estarem relacionados à técnica. De forma geral, o modelo de \gls{ttt} dos artigos revisados é extremamente semelhante, variando nas aplicações, algumas considerações, e outras pequenas peculiariedades.

\subsection{Breve descrição dos artigos revisados}
\label{sec_ttt:subsec_revisao}
Como discutido na \cref{subsec:sec1:estimacao_ssf_por_inversao}, a técnica de \gls{ttt} é versátil devido à sua simplicidade, sendo aplicável em uma grande gama de aplicações, de sismologia acústica a tomografia computadorizada de tecidos, e avaliação eletromagnética de propriedades da água. Em específico para o caso de meios sólidos com ondas acústicas, a \gls{ttt} utiliza primariamente ondas P, dado que essas viajam mais rapidamente, e essa técnica baseia-se no tempo de trânsito entre fontes e receptores. Abaixo segue uma breve descrição dos artigos selecionados.

\begin{itemize}[leftmargin=4pt+\widthof{\cite{aki_determination_1976,aki_determination_1977}}]
	\citem{aki_determination_1976,aki_determination_1977} - Esses dois trabalhos são considerados os artigos pioneiros da \gls{ttt} (sendo publicados respectivamente em \citeyear{aki_determination_1976} e \citeyear{aki_determination_1977}), descrevendo a classificação de camadas geológicas por meio da estimação do \gls{ssf}, e a estimação da localização do epicentro de terremotos.
	\citem{ali_opensource_2019} - Descreve a estimativa do SSF em tecidos utilizando tomografia computadorizada com ultrassom, por meio de \gls{fmm} para ray-tracing e um solver eikonal.
	\citem{hormati_robust_2010} - Trata de um problema semelhante a \cite{ali_opensource_2019}, descrevendo diferentes abordagens para resolver o problema inverso.
	\citem{jeong_investigating_2024} - Compara \gls{ttt} com tomografia por reflexão, usando ambas para treinamento de uma CNN.   
	\citem{hellmann_borehole_2023} - Descreve o uso de \gls{ttt} para estimação do \gls{ssf} em furos de sondagem em geleiras, com um passo de correção da posição dos dispositivos ao longo do furo.
	\citem{dines_computerized_1979} - Descreve o uso de  \gls{ttt}, para estimar o campo de atenuação das ondas eletromagnéticas, além do campo de velocidade. Considera sinais eletromagnéticos, mas a teoria e a matemática são as mesmas. Também propõe uma solução iterativa em pares de dispositivos --- ou seja, fonte-receptor, eles iterativamente atualizam e refinam o campo, ao invés de utilizar toda a informação simultaneamente.
	\citem{phillips_traveltime_1991} - Relata  diferentes métodos e modificações na \gls{ttt} de forma geral, comparando seus resultados em métricas como rugosidade, erro, espalhamento; com cada método trocando desempenho global por melhoria em alguma métrica.
	\citem{tang_travel_2024} - Propõe o uso da técnica \gls{fsm} ao invés de \gls{fmm} para solução da equação eikonal no problema direto, e foca primariamente na comparação dos dois solvers. Modelagem de reflexões com o FSM também é discutida no artigo.
	\citem{zhang_inversion_2016} - Propõe outra adaptação da \gls{ttt} para inclusão de reflexões. O modelo minimiza uma combinação de vagarosidades média e aparente, além do tempo de regularização. Parte do princípio que, em certas aplicações, ondas refletidas podem chegar em momentos próximos que as diretas, permitindo seu uso.
\end{itemize}

\subsubsection*{Temas em comum}
%
\def\X{\ensuremath{\mathsf{X}}}
%
A \cref{table:fttt_revisao_temas} compila uma lista dos temas em comum apresentados nos diferentes artigos, e um pequeno comentário sobre os aspectos ou considerações que mais destacam cada. Note que em todos os artigos apresentados, assume-se a discretização do espaço. Nem todos consideram efeitos de curvatura dos caminhos das ondas (\ingles{ray-bending}), dependendo da aplicação; alguns realizam minimização iterativa através do gradiente conjugado, enquanto outros buscam uma formulação fechada para o problema inverso. A maioria utiliza alguma forma de regularização na minimização, e alguns também propõem o uso de \ingles{blur}, dispersando os raios estimados.

\begin{table}[!t]\small\centering
	\caption{Tabela comparativa dos artigos analisados. CG - Gradiente Conjugado (para minimização); $P_{\alpha}$ - Regularização; $\bvB{i}$ - Suavização do \gls{ssf}.}
	\label{table:fttt_revisao_temas}
	\begin{tabular}{l c c c c c l}
		Paper 								& Discreto 	& \textsl{Ray bending} 	& CG 	& $P_{\alpha}$ 	& $\bvB{i}$ & Extra 	\\
		\cite{aki_determination_1976} 		& \X 		& -- 			& --	& \X			& --		& --		\\
		\cite{aki_determination_1977} 		& \X 		& -- 			& --	& \X			& --		& Análise de perturbação		\\
		\cite{ali_opensource_2019}			& \X		& \X			& \X	& \X			& \X		& Linearização estatística \\
		\cite{hormati_robust_2010}			& \X 		& \X			& \X	& \X			& --		& --		\\
		\cite{jeong_investigating_2024}	& \X		& \X			& --	& \X			& --		& --		\\
		\cite{hellmann_borehole_2023}		& \X		& \X			& --	& \X			& --		& Estimação de posição	\\
		\cite{dines_computerized_1979}		& \X		& --			& --	& --			& --		& Considera atenuação\\
		\cite{phillips_traveltime_1991}	& \X		& --			& --	& \X			& \X		& Compara várias técnicas	\\
		\cite{tang_travel_2024}			& \X		& \X			& \X	& --			& --		& Usa FSM \\
		\cite{zhang_inversion_2016}		& \X		& \X			& \X	& \X			& --		& Considera reflexões			
	\end{tabular}
\end{table}

\subsection{Fundamentos da \textsl{Travel-Time Tomography}}
Na \gls{ttt}, o ambiente é divido em uma grade de células (também aplicável para 3D ou 4D; e também para casos 1D, em que se divide o ambiente em camadas), e uma {distribuição de dispositivos} (fontes e receptores) dentro do ambiente (geralmente no entorno dele). Assume-se que o SSF é portanto {discretizado espacialmente}, com a velocidade de cada célula sendo constante.

Cada fonte produz um sinal sonoro, que é captado pelos receptores\footnote{A mecânica de produção, captação, ou demultiplexação dos sinais não entra no escopo da \gls{ttt}.}, onde mensura-se o \gls{tempo de trânsito} (ou tempo de primeira chegada) entre a fonte e o receptor; esse tempo de trânsito é a informação primária utilizada na estimação do \gls{ssf} dentro da \gls{ttt}.

A modelagem da \gls{ttt} parte do {princípio de Fermat de tempo mínimo}, em que o tempo de primeira chegada da onda no trajeto fonte-receptor será o caminho direto (mas curvado) entre os dois dispositivos. Para o funcionamento dessa técnica, também assume-se a ausência de reflexões (internas ou externas) das ondas.

\subsubsection*{Problema direto}

O processo de estimação do \gls{ssf} pela \gls{ttt} pode ser separado em duas etapas: o {problema direto}, em que estima-se os tempos de trânsito entre as fontes e receptores, dado uma estimativa do SSF; e o {problema inverso}, em que é construído um \gls{ssf} baseado nos tempo de trânsito (e nas distâncias percorridas por cada raio em cada célula). Essas duas etapas são repetidas iteradamente, utilizando a solução anterior de uma como os dados de entrada da outra, até alcançar convergência.

No problema direto, é necessário algum solver da equação eikonal, essa sendo a equação que rege a relação entre o \gls{ssf} e os tempos de primeira chegada considerando também curvas. Na literatura, é amplamente utilizado o solver eikonal \gls{fmm}, que baseia-se em uma otimização por grafos; e ocasionalmente o \gls{fsm}. Raramente também surgem outros artigos com solvers da equação eikonal autorais, porém esses são geralmente desenhados para aplicações específicas.

\subsubsection*{Problema inverso}

Como já descrito, o problema inverso objetiva estimar o \gls{ssf} utilizando: a solução obtida para a equação eikonal, os tempos de trânsito medidos, e as simplificações provenientes da discretização espacial do ambiente. Também é comum encontrar na literatura termos de regularização (de forma a condicionar a solução obtida, trabalhando alguma métrica secundária) e por vezes algum método de \ingles{blur} (também uma alternativa para melhorar a qualidade dos dados de entrada).

Na literatura, há dois caminhos que podem ser seguidos para encontrar a solução desse problema inverso: embora a maioria utilize formas fechadas para obtenção da solução em termos das variáveis conhecidas (solução da equação eikonal, e tempos de trânsito medidos) através de uma modelagem quadrática do erro e minimização pelo Jacobiano; algumas seguem um caminho iterativo dentro do problema inverso, geralmente motivadas por alguma outra informação adicional considerada. Contudo, a maioria dos problemas que utilizam a solução iterativa podem ser retrabalhados de forma a encontrar uma solução fechada.

\subsection{\textsl{Mindmap}}

A \cref{fig:mindmap_sec2} apresenta um \textsl{mindmap} dos conceitos expostos até o momento na \cref{sec:travel-time_tomography}, dispondo graficamente a organização das informações apresentadas. A cor de cada nó indica o conceito-pai dele. Note que alguns nós (em específico: tempo de trânsito, solução eikonal, e estimação do \gls{ssf}) estão ligados a mais um nó; nesses casos a cor ainda indica o nó-pai, e a conexão indica que o outro nó está intrinsecamente ligado a esse conceito.

%\section{Mindmap}
%\vfill
\begin{figure}[!t]
	\centering
	\begin{tikzpicture}[
			mindmap,
			text width=2.5cm,
			align=center,
			every node/.style={concept},
			concept color=green!20!white!85!black,
			level 1/.append style={
				level distance=5cm,
				minimum size=3cm,
				text width=2.8cm,
				font=\large,
			},
			level 2/.append style={
				level distance=3.5cm,
				minimum size=2cm,
				text width=1.8cm
			},
			]
			
			\node[minimum size=4.5cm, text width=4.3cm, line width=6pt, fill=green!10!white!95!black] (TTT) {\Large\textbf{\textsl{Travel-Time \mbox{Tomography} (TTT)}}}
			child[concept color=yellow!10!white!95!black,grow=45] { 
				node (fundamentos) {\textbf{Fundamentos}}
				child[grow=150] { node{Distribuição\\de dispositivos} }
				child[grow=30] { node{Discretização\\espacial} }
				child[grow=-30] { node{Princípio\\de Fermat} }
				child[grow=-90] { node (tt) at (fundamentos|-TTT) {Tempo de\\trânsito} }
			}
			child[concept color=teal!10,grow=315] { 
				node (probinv) {\textbf{Problema\\inverso}}				
				child[grow=-30] { node{Otimização\\quadrática} }
				child[grow=30] { node{Regularização} }
				child[concept color=yellow!10!white!95!black] { node at (tt) {Tempo de trânsito}}
				child[concept color=orange!10] { node (solucaoeikonal) at (TTT|-probinv) {Solução\\eikonal} }
				child[grow=225] { node (estimssf) {Estimação\\do SSF} }
			}
			child[concept color=orange!10,grow=225] { 
				node (probdir) {\textbf{Problema\\direto}} 
				child[grow=300] { node{\textsl{Ray-tracing}} }
				child[grow=240] { node{FMM/FSM} }
				child { node at (solucaoeikonal) {Solução\\eikonal} }
				child[concept color=teal!10] { node at (estimssf) {Estimação\\do SSF} }
%				child[concept color=teal!10,grow=315, level distance=5cm] { node(estssf) {} }
			}
			;	
	\end{tikzpicture}
	\caption[\textsl{Mindmap} da \cref{sec:travel-time_tomography}.]{\textsl{Mindmap} da \cref{sec:travel-time_tomography} sobre a tomografia de tempo de trânsito.}
	\label{fig:mindmap_sec2}
\end{figure}

